\documentclass[a4paper]{article}     
\usepackage{amsfonts}
\usepackage{amsmath}
\usepackage{amssymb}
\usepackage{graphicx}
\usepackage{cancel}

\newcommand{\asa}{\Leftrightarrow} %als slechts als
\newcommand{\adR}{\Rightarrow} %als dan Rechts
\newcommand{\adL}{\Leftarrow}
\newcommand{\Lh}{\left(} %Links haakje
\newcommand{\Rh}{\right)}
\newcommand{\Lv}{\left[} %Links vierkanthaakje
\newcommand{\Rv}{\right]}
\newcommand{\La}{\left\{} %Links acolade
\newcommand{\Ra}{\right\}}
\newcommand{\Ras}{\right.} %Rechts acolade stelsel (omdat om stelsels te maken heb je een punt nodig
\newcommand{\pnR}{\rightarrow} %pijl naar Rechts
\newcommand{\limiet}{\lim\limits_}
\newcommand{\schrap}{\cancel}
\newcommand{\schrapb}{\bcancel} %backslash
\newcommand{\schrapk}{\xcancel} %kruis
\newcommand{\vervang}{\cancelto} 

\begin{document}

\noindent \large Auteur: Yoren Collier \\
\noindent \large Studierichting: Informatica\\
\noindent \large Oefeningenreeks 3B nr. 3\\

\medskip

\normalsize

$f(x) = \dfrac{1}{\sqrt{x}} = x^{\tfrac{-1}{2}}$\\

$\dfrac{df}{dx}(x) = \dfrac{-1}{2} \Lh x \Rh ^{\tfrac{-3}{2}}$\\

$\dfrac{d^2f}{dx^2}(x) = \dfrac{-1}{2} \cdot \dfrac{-3}{2} \Lh x \Rh ^{\tfrac{-5}{2}}$\\

$\dfrac{d^3f}{dx^3}(x) = \dfrac{-1}{2} \cdot \dfrac{-3}{2} \cdot \dfrac{-5}{2} \Lh x \Rh ^{\tfrac{-7}{2}}$\\

$\dfrac{d^4f}{dx^4}(x) = \dfrac{-1}{2} \cdot \dfrac{-3}{2} \cdot \dfrac{-5}{2} \cdot \dfrac{-7}{2} \Lh x \Rh ^{\tfrac{-9}{2}} = \dfrac{+105}{16 \sqrt{x^9}} = \dfrac{105}{16x^4 \sqrt{x}} = \dfrac{105 \sqrt{x}}{16x^5}$\\

\begin{thebibliography}{999}
%\bibitem{a} Referentie
\end{thebibliography}
\end{document} 